\documentclass{article}
\usepackage{../../Self_Style}

\title{Phys 115B HW 2}
\author{Zih-Yu Hsieh}
\date{\today}

\begin{document}
\maketitle

\section*{1}

\begin{ques}\label{q1}
More Practice with Spherical Harmonics

Consider a particle in three dimensions in a state
\[
\psi(x,y,z)=A(x+y+2z)e^{-\alpha r}
\]
where $\alpha>0$ is a positive constant, and $A$ a normalization factor.

\begin{itemize}
\item[(a)] Write the spherical harmonics $Y^{0,\pm1}_{1}$ in terms of $x,y,z$, and
\[
r=\sqrt{x^{2}+y^{2}+z^{2}}.
\]

\item[(b)] Using your result from part (a), show that $\psi$ can be written as a product of a purely
$r$-dependent wave function and a linear combination of spherical harmonics.
\end{itemize}
\end{ques}

\begin{proof}
\end{proof}

\newpage

\begin{ques}\label{q2}
Radial Momentum

In lecture, we defined an operator corresponding to the momentum in the radial direction,
\[
p_r \equiv \frac{1}{2}(\hat{r}\cdot p + p\cdot \hat{r})
= -\frac{i\hbar}{2}\left(\frac{1}{r}\, r\cdot\nabla + \nabla\cdot\frac{r}{r}\right).
\]

\begin{itemize}
\item[(a)] Show that $\hat{r}\cdot p$ is not Hermitian, and hence cannot itself be identified with
momentum in the radial direction.

\item[(b)] Using the definition of $p_r$ above, show that
\[
p_r^2 = -\hbar^2 \frac{1}{r^2}\frac{\partial}{\partial r}
\left(r^2\frac{\partial}{\partial r}\right).
\]
\end{itemize}
\end{ques}

\begin{proof}
\end{proof}

\newpage

\begin{ques}\label{q3}
The Size of Hydrogen (Griffiths 4.15, modified, 4.13 in second edition, still modified). Do parts (a)--(c) from the Griffiths problem.

\begin{itemize}
\item[(a)] Find $\left<r\right>$ and $\left<r^2\right>$ for an electron in the ground state of hydrogen. Express your answers in terms of te Bohr radius.
\item[(b)] Find $\left<x\right>$ and $\left<x^2\right>$ for an electron in the ground state of hydrogen. 

(\emph{Hint}: This requires no new integration - note that $r^2=x^2+y^2+z^2$, and exploit the symmetry of the ground state).

\item[(c)] Find $\left<x^2\right>$ in the state $n=2,\ell=1,m=1$. 

(\emph{Hint:} this staet is \emph{not} symmetrical in $x,y,z$. Use $x=r\sin(\theta)\cos(\phi)$.

\item[(d)] Find $\langle r\rangle$ and $\langle r^2\rangle$ for an electron in a circular orbit of hydrogen with arbitrary principal quantum number $n$ (corresponds to $l=n-1$ and any allowed $m$).

\item[(e)] Compute the RMS uncertainty
\[
\sqrt{\langle r^2\rangle - \langle r\rangle^2}
\]
in terms of $r$ for the electron in part (d). Note that the fractional spread in $r$ decreases with increasing $n$ (in this sense the system ``begins to look classical'' for large $n$). How much more volume does a hydrogen atom in the $n=100$ state occupy compared to the hydrogen atom in the ground state. (Hint --- you might want to look at Griffiths 4.55, or 4.15 in the second edition.)
\end{itemize}
\end{ques}

\begin{proof}
\end{proof}

\newpage

\begin{ques}\label{q4}
The Size of Hydrogen continued (Griffiths 4.16, 4.14 in the second edition)

What is the most probable value of $r$, in the ground state of hydrogen? (The answer is not
zero!) Hint: First you must figure out the probability that the electron would be found
between $r$ and $r+dr$.
\end{ques}

\begin{proof}
\end{proof}

\newpage

\begin{ques}\label{q5}
Where is the Electron (Griffiths 4.54, 4.45 in the second edition)

What is the probability that an electron in the ground state of hydrogen will be found \emph{inside the nucleus}?
\begin{itemize}
    \item[(a)] First calculate the \emph{exact answer}, assuming the wave function is correct all the way down to $r=0$. Let $b$ be the radius of the nucleus.
    
    Here is the wave function:
    $$\psi_{1,0,0}(r,\theta,\phi) = \frac{1}{\sqrt{\pi a^3}}e^{-r/a}$$ 
    \item[(b)] Expand your result as a power series in the small number $\epsilon := \frac{2b}{a}$, and show that the lowest-order term is the cubic: $p \approx \frac{4}{3}\left(\frac{b}{a}\right)^3$. This should be a suitable approximation, provided that $b\ll a$ (which it is).
    \item[(c)] Alternatively, we mmight assume that $\psi(r)$ is essentially constant over the (tiny) volume of the nucleus, so that $P\approx \frac{4}{3}\pi b^3|\psi(0)|^2$. Check that you get the same answer this way.
    \item[(d)] Use $b\approx 10^{-15}$ m and $a\approx 0.5\cdot 10^{-10}$ m to get a numerical estimate for $P$. Roughly speaking, this represents the ``fraction of its time that the electron spends inside the nucleus''.   
\end{itemize}
\end{ques}

\begin{proof}
\end{proof}

\end{document}

